problem:  
Let \( M^n \) be a compact, simply connected Riemannian manifold with positive sectional curvature that admits an isometric cohomogeneity‑one action by a compact Lie group \( G \).  For known examples (rank‑one symmetric spaces, certain Eschenburg and Bazaikin biquotients), the isotropy subgroups \( (H \subset K_\pm \subset G) \) have highly restricted representation types, often forcing \( G \) to have rank at most two.  

**Question:**  
Does positive sectional curvature in cohomogeneity one force an *upper bound on the rank of the acting group \(G\)* that is *strictly smaller* than the general Grove–Searle symmetry‑rank bound?  
Equivalently, does positive curvature prohibit high‑rank isotropy representations from appearing in the group diagram \( (G,H,K_+,K_-) \)?  
More concretely, is it true that for sufficiently large \(n\), a positively curved cohomogeneity‑one \(n\)‑manifold must be diffeomorphic to a rank‑one symmetric space or a known biquotient whose \(G\)-action has \(\mathrm{rank}(G)\le 2\)?  

Why is it a "good" problem:  
This formulation shifts the focus from a purely numerical bound to a **structural obstruction** at the level of isotropy representations, bridging transformation group theory, representation theory, and Riemannian geometry.  If such a rank bound exists, it would reveal a deep rigidity principle: large symmetry groups cannot act with low cohomogeneity under positive curvature because the isotropy representations would be unable to sustain curvature positivity.  Conversely, if such a restriction fails, any counterexample would yield genuinely new families of positively curved manifolds with unprecedented symmetry structures.  Thus the problem probes the hidden compatibility conditions among curvature, symmetry, and orbit geometry—an uncharted connection that could either expose a new rigidity theorem or lead to the discovery of new manifolds.  It is concise, conceptually rich, and represents a clear frontier question at the core of global differential geometry.